
\chapter{Part VIII Review: Beyond Perturbative Gauge Theory}

The final part adds three ideas that cannot be reduced to a single Feynman rule: gauge redundancy in non-Abelian theories, symmetry breaking, and nonperturbative sectors.

\section{A guided reconstruction}

% QFTFIGURE-BEGIN partreview08-01-part-map
\qftfigure{figures/instructional/partreview08/partreview08-01-part-map.pdf}{The map collects the part's main constructions into a visual route so the reader can see how the chapters support one another.}{fig:partreview08-01-part-map}
% QFTFIGURE-END partreview08-01-part-map












The Yang-Mills field strength contains a commutator, so gauge bosons self-interact.  Gauge fixing introduces ghost fields in covariant gauges.  Spontaneous symmetry breaking reorganizes degrees of freedom around a chosen vacuum.  Nonperturbative physics then asks what is missed by expanding around only one saddle.

\begin{exercise}
Explain why instanton-like contributions are invisible in a power series in \(g\).
\begin{answer}
They scale like \(\exp(-\hbox{constant}/g^2)\), which cannot be reproduced by any finite Taylor series in powers of \(g\) around \(g=0\).
\end{answer}
\end{exercise}
